\chapter{Технологический раздел}

В данном разделе представляются выбор средств реализации, тестирование и демонстрация работы программного обеспечения, а также результаты обучения нейронной сети.

\section{Выбор средств программной реализации}

В качестве языка программирования для реализации программного обеспечения был выбран Python~\cite{python} по следующим причинам:
\begin{itemize}
	\item поддержка объектно-ориентированной методологии проектирования;
	\item наличие в библиотеке языка PyTorch~\cite{pytorch} необходимых классов предобученных искусственных нейронных сетей, выбранных по результатам проектирования;
	\item наличие в библиотеке языка OpenCV~\cite{opencv} необходимых функций реализации алгоритма Виолы-Джонса, выбранного по результатам проектирования.
\end{itemize}

\section{Результаты обучения нейронной сети}

Качество обучения нейронной сети для решения задачи бинарной классификации характеризуется набором количественных метрик, описываемых с помощью следующих терминов:
\begin{itemize}
	\item \textit{истинно положительное срабатывание} -- событие, когда возвращаемое нейронной сетью значение является положительным и соответствует действительности;
	\item \textit{истинно отрицательное срабатывание} -- событие, когда возвращаемое нейронной сетью значение является отрицательным и соответствует действительности;  
	\item \textit{ложно положительное срабатывание} -- событие, когда возвращаемое нейронной сетью значение является положительным и не соответствует действительности; 
	\item \textit{ложно отрицательное срабатывание} -- событие, когда возвращаемое нейронной сетью значение является отрицательным и не соответствует действительности.
\end{itemize}

\textit{Точностью} искусственной нейронной сети называется следующая метрика~\cite{ml}:
\begin{equation}
	Accuracy = \frac{TP + TN}{TP + TN + FP + FN},
\end{equation}
где
\begin{itemize}
	\item $TP$ -- количество истинно положительных срабатываний;
	\item $TN$ -- количество ложно положительных срабатываний;
	\item $FP$ -- количество истинно отрицательных срабатываний;
	\item $FN$ -- количество ложно отрицательных срабатываний.
\end{itemize}

\textit{Полнотой} искусственной нейронной сети называется следующая метрика~\cite{ml}:
\begin{equation}
	Recall = \frac{TP}{TP + FN},
\end{equation}
где
\begin{itemize}
	\item $TP$ -- количество истинно положительных срабатываний;
	\item $FN$ -- количество ложно отрицательных срабатываний.
\end{itemize}
Данная метрика отражает долю верно определенных образцов положительного класса среди всех образцов положительного класса.

\textit{Сходимостью} искусственной нейронной сети называется следующая метрика~\cite{ml}:
\begin{equation}
	Precision = \frac{TP}{TP + FP},
\end{equation}
где
\begin{itemize}
	\item $TP$ -- количество истинно положительных срабатываний;
	\item $FP$ -- количество истинно отрицательных срабатываний.
\end{itemize}
Данная метрика отражает долю  верно определенных образцов положительного класса среди всех образцов, отнесенных к этому классу.

\textit{F-мерой} искусственной нейронной сети называется метрика равная среднему гармоническому между её полнотой и сходимостью~\cite{ml}:
\begin{equation}
	\textit{F-score} = 2 \cdot \frac{Precision \cdot Recall}{Precision + Recall},
\end{equation}
\begin{itemize}
	\item $Precision$ -- сходимость ИНС;
	\item $Recall$ -- полнота ИНС.
\end{itemize}
Данная метрика отражает способность искусственной нейронной сети достигать высоких полноты и сходимости.

Зависимости значений метрик обучения разработанной нейронной сети от эпохи представлены в таблицах~\ref{tab:training-history-full}~и~\ref{tab:training-history}.

\begin{table}[h]
	\centering
	\caption{Зависимости значений метрик качества обучения разработанной нейронной сети от эпохи}
	\label{tab:training-history}
	\begin{tabularx}{\textwidth}{|c|Y|Y|Y|Y|}
		\hline
		\textbf{Эпоха} & \textbf{Потеря на тренировочной выборке} & \textbf{Потеря на валидационной выборке} & \textbf{Точность на тренировочной выборке, \%} & \textbf{Точность на валидационной выборке, \%} \\
		\hline
		5  & 0.93616 & 0.93795 & 68.97 & 68.87 \\
		\hline
		10 & 0.72180 & 0.72417 & 75.88 & 75.61 \\
		\hline
		15 & 0.54823 & 0.55281 & 80.97 & 80.76 \\
		\hline
		20 & 0.42918 & 0.42835 & 85.19 & 85.02 \\
		\hline
		25 & 0.33418 & 0.33477 & 88.73 & 88.37 \\
		\hline
		30 & 0.25464 & 0.25741 & 91.12 & 90.97 \\
		\hline
		35 & 0.19737 & 0.20258 & 93.16 & 93.02 \\
		\hline
		40 & 0.15842 & 0.16019 & 94.48 & 94.31 \\
		\hline
		45 & 0.12447 & 0.12309 & 95.75 & 95.51 \\
		\hline
		50 & 0.10512 & 0.09613 & 96.92 & 96.57 \\
		\hline
		55 & 0.08707 & 0.09144 & 97.31 & 97.12 \\
		\hline
		60 & 0.07671 & 0.07164 & 98.05 & 98.67 \\
		\hline
		65 & 0.05816 & 0.05242 & 98.36 & 98.54 \\
		\hline
	\end{tabularx}
\end{table}

По таблице~\ref{tab:training-history-full} были построены графики~\ref{img:loss}~и~\ref{img:accuracy} зависимостей значений функции потерь и значений точности от эпохи соответственно.

\includeimage
{loss}
{f} % Обтекание (без обтекания)
{H}
{\textwidth} % Ширина рисунка
{График зависимости значений функции потерь от эпохи}

\includeimage
{accuracy}
{f} % Обтекание (без обтекания)
{H}
{\textwidth} % Ширина рисунка
{График зависимости значений точности от эпохи}

Результаты обучения разработанной нейронной сети представлены в таблицах~\ref{tab:training-results-1}~и~\ref{tab:training-results-2}.

\begin{table}[h]
	\centering
	\caption{Результаты обучения разработанной нейронной сети}
	\label{tab:training-results-1}
	\begin{tabularx}{\textwidth}{|Y|Y|Y|Y|}
		\hline
		\textbf{Точность, \%} & \textbf{Полнота, \%} & \textbf{Сходимость, \%} & \textbf{F-мера, \%} \\
		\hline
		98.67 & 98.65 & 98.65 & 98.65 \\
		\hline
	\end{tabularx}
\end{table}

\begin{table}[h]
	\centering
	\caption{Результаты обучения разработанной нейронной сети по классам}
	\label{tab:training-results-2}
	\begin{tabularx}{\textwidth}{|Y|Y|Y|Y|}
		\hline
		\textbf{Класс} & \textbf{Полнота, \%} & \textbf{Сходимость, \%} & \textbf{F-мера, \%} \\
		\hline
		Подлинные изображения & 98.50 & 98.79 & 98.65 \\
		\hline 
		Синтезированные изображения & 98.79 & 98.50 & 98.65 \\
		\hline
	\end{tabularx}
\end{table}

Из полученных результатов, можно сделать вывод, что обучение модели характеризуется монотонным снижением значений функции потерь на тренировочной и валидационной выборках на всех эпохах без признаков переобучения. Достигнутые значения полноты и сходимости для обоих классов близки, что подтверждает сбалансированность классификации без смещения в сторону одного из типов изображений. Разработанная нейронная сеть отвечает предъявленным требованиям к точности и к вычислительной сложности, описанным в разделе~\ref{sec:method-requirements}.

\section{Демонстрация работы программного обеспечения}

Демонстрация работы разработанного программного обеспечения представлена на рисунке~\ref{img:mainwindow}~и~\ref{img:examples}.

\includeimage
{mainwindow}
{f} % Обтекание (без обтекания)
{H}
{0.95\textwidth} % Ширина рисунка
{Главное окно приложения}

Главное окно приложения позволяет пользователю загрузить изображение для обнаружения признаков генерации и впоследствии демонстрирует ему результаты проверки как всего изображения в целом, так и отдельных обнаруженных на нем лиц. Панель управления дает возможность пользователю изменить значение пороговой вероятности функции дискриминации для формирования результирующего решения о поддельности изображения. 

\includeimage
{examples}
{f} % Обтекание (без обтекания)
{H}
{\textwidth} % Ширина рисунка
{Результаты обнаружения признаков генерации на изображениях с использованием разработанного программного обеспечения при значении пороговой вероятности равном 0.5}

Из полученных результатов можно сделать вывод, что разработанное программное обеспечение соответствует спроектированной архитектуре и отвечает всем предъявленным требованиям, описанным в разделе~\ref{sec:software-requirements}. Исходный код приложения доступен в открытом сетевом репозитории~\cite{git}.

\section{Тестирование программного обеспечения}

Были разработаны модульные и интеграционные тесты для методов классов, реализующих компоненты доступа к данным и бизнес-логики. Модульное тестирование выполнялось в соответствии с \textit{лондонской моделью}~\cite{tests}, а интеграционное -- с взаимодействием с файловой системой.

Сквозное тестирование разработанного программного обеспечения производилось в соответствии с сформированными классами эквивалентности, приведенными в таблице~\ref{tab:equal-classes}. 

\begin{table}[h]
	\centering
	\caption{Классы эквивалентности сквозного тестирования разработанного программного обеспечения}
	\label{tab:equal-classes}
	\begin{tabularx}{\textwidth}{|Y|Y|}
		\hline
		\textbf{Класс} & \textbf{Описание} \\
		\hline
		Подлинные изображения & Изображения корректного формата, содержащие одно или несколько лиц и не имеющие признаки генерации \\
		\hline
		Поддельные изображения & Изображения корректного формата, содержащие одно или несколько лиц и имеющие признаки генерации \\
		\hline 
		Изображения без лиц & Изображения корректного формата, не содержащие лица \\
		\hline
		Некорректный файл & Файлы с расширением или содержимым, отличным от ожидаемого \\
		\hline
		Граничные значения пороговой вероятности & Установка пороговой вероятности равной граничным значениям \\
		\hline
	\end{tabularx}
\end{table}

Все проведенные тесты были успешно пройдены. Из результатов тестирования разработанного программного обеспечения можно сделать вывод, что оно корректно обрабатывает все выделенные классы эквивалентности. Успешное прохождение модульных и интеграционных тестов подтверждает надежность компонентов доступа к данным и бизнес-логики. 

\section{Выводы из технологического раздела}

В результате реализации метода обнаружения признаков генерации на изображении с использованием нейронной сети были решены следующие задачи:
\begin{itemize}
	\item реализована и обучена необходимая искусственная нейронная сеть;
	\item разработано и протестировано программное обеспечение, реализующее предложенный метод.
\end{itemize}

Результаты обучения нейронной сети и тестирование разработанного программного обеспечения показали соответствие метода всем предъявленным требованиям и его практическую применимость.